\section{Background}
\label{sec:background}

\paragraph{The update.} Let $\theta$ denote the model's $P$ parameters and $f$ the loss to
be minimized, which ES calls the fitness. Each generation evaluates the fitness of $N$
candidates, $f_i=f(\theta+\sigma\epsilon_i)$, where $\epsilon_i$ is candidate $i$'s random
perturbation, an array with the shape of $\theta$, and $\sigma$ sets the noise scale. It
then changes the parameters by
\begin{equation}
  \Delta\theta = -\frac{\eta}{N\sigma}\sum_{i=1}^{N} w_i\,\epsilon_i ,
  \label{eq:update}
\end{equation}
where $\eta$ is the learning rate and $w_i$ is a coefficient computed from the fitness
scores \citep{wierstra2014,nesterov2017}. A reward to be maximized enters negated. We use
\emph{centered ranks}: each score is replaced by its rank among the $N$ scores, scaled
linearly to $[-\tfrac12,\tfrac12]$ \citep{salimans2017}. The update therefore moves toward
perturbations with lower losses and away from those with higher losses. Because a rank
depends on every score, all $N$ scores must reach every device before any coefficient can
be computed, so both placements begin by sharing them.

\emph{Mirrored} sampling evaluates each perturbation $\epsilon$ together with
$-\epsilon$, so $N$ candidates use $N/2$ independent directions. The terms for $\epsilon$
and $-\epsilon$ in Eq.~\ref{eq:update} combine into one, weighted by the difference of their
coefficients, so reconstruction processes $N/2$ perturbations instead of $N$.

\paragraph{Two separate choices.} The perturbation design determines how perturbations
are generated, stored and evaluated, and so how much work reconstruction takes. The
update placement determines where the weighted sum in Eq.~\ref{eq:update} is computed:
in full on every device (replication), or in parts that an all-reduce adds (splitting).
We vary the two choices independently.

Full-rank Gaussian noise perturbs every weight independently. During evaluation it can
be held as a full array for every candidate (\emph{dense}) or regenerated from a seed as
each candidate is evaluated (\emph{seed}) \citep{qiu2025}. Seed needs memory for only one
candidate's noise at a time, but evaluates the candidates one at a time. Both variants
generate the noise a second time, from seeds, for reconstruction.

Low-rank noise represents the perturbation of an $m\times n$ weight matrix as
$UV^\top/\sqrt r$, where $U$ ($m\times r$) and $V$ ($n\times r$) have independent
standard Gaussian entries; the factor $1/\sqrt r$ gives each entry of the product unit
variance, as in full-rank noise. The factors store $r(m+n)$ values instead of $mn$, and
the full perturbation matrix is never formed \citep{eggroll2025}. Each perturbation has
rank at most $r$, but the weighted sum over candidates can have higher rank, so the
update is still a full $m\times n$ matrix. Table~\ref{tab:arms} lists the names used
below; unless stated otherwise, \emph{rank $r$} refers to mirrored sampling.

\begin{table}[t]
\centering
\small
\setlength{\tabcolsep}{3pt}
\caption{Perturbation variants. Dense and seed draw the same full-rank noise; during
  evaluation, dense holds every candidate's noise array at once and seed holds one at a
  time. Every variant regenerates its noise from seeds for reconstruction. Mirrored
  sampling is chosen independently.}
\label{tab:arms}
\begin{tabular}{llll}
\toprule
name & noise & held during & mirrored? \\
 & & evaluation & \\
\midrule
dense & full rank & all arrays & no \\
seed & full rank & one array & no \\
seed, mirrored & full rank & one array & yes \\
rank $r$ & low rank & thin factors & yes \\
rank $r$, non-mirrored & low rank & thin factors & no \\
\bottomrule
\end{tabular}
\end{table}
